\section{Step 28 Mix the first half of the diagonal groups}\label{chacha-s28}

Run the same first two vector groups after selection. The arithmetic code is unchanged, but lane zero now represents a diagonal rather than a column. Preserve the shuffled layout until all four groups complete.

\subsection{Write and inspect}
Replace the body of region \texttt{msa\_diagonal\_a}. Retain its \texttt{\# BEGIN} and \texttt{\# END} comments and the supplied labels outside them.

\Needspace{15\baselineskip}
\MARSCodeLabel{MIPS assembly}
\begin{minted}[samepage=true]{mars-idea}
    addv.w $w0, $w0, $w1
    xor.v $w3, $w3, $w0
    slli.w $w8, $w3, 16
    srli.w $w9, $w3, 16
    or.v $w3, $w8, $w9
    addv.w $w2, $w2, $w3
    xor.v $w1, $w1, $w2
    slli.w $w8, $w1, 12
    srli.w $w9, $w1, 20
    or.v $w1, $w8, $w9
\end{minted}

Use recipe \texttt{s28-1} with the cumulative program through this edit. Inspect its named checkpoint before continuing to the temporary or normal return.

\subsection{Expected observation}
All four diagonal selections have completed their first two groups.

\Needspace{9\baselineskip}
\subsection{Check your understanding}
Why can the arithmetic sequence be reused?

\begin{itemize}
\item \textbf{A} Columns and diagonals have identical original values.
\item \textbf{B} The shuffle places the required four words in corresponding lanes.
\item \textbf{C} The shuffle computes the additions.
\end{itemize}
